El acceso abierto a recursos científicos requiere organización para sostener la participación. Este ensayo compara CoAfina, el editatón Salud para todos con la OMS y Cornell, y experiencias venezolanas con ATLAS Open Data. Una síntesis documental distingue inscripción, actividad, productos y continuidad. Recuperamos resultados públicos de tres ediciones de CoAfina, definiciones y metas publicadas del editatón, y cocientes descriptivos con sus límites. Los editatones registran cocientes de actividad del 42,2\,\% y 26,5\,\%; los umbrales de otros formatos difieren e impiden comparar eficacia educativa. La documentación identifica ayudas asociadas a entregas, un canal comunitario que persistió y un paso de participante a proponente de retos. Una matriz relaciona estos mecanismos con preservación, reconocimiento y seguimiento. Proponemos reutilizar definiciones operativas, vincular cada producto a sus materiales y asignar responsabilidades posteriores. El marco ofrece decisiones adaptables a instituciones venezolanas sin confundir alcance, aprendizaje y continuidad.
